\documentclass[a4paper,10pt]{amsart}
\pdfinfoomitdate=1
\pdftrailerid{}
\pdfsuppressptexinfo=15
\usepackage[T1]{fontenc}
\usepackage{lmodern}
\usepackage[margin=25mm]{geometry}
\usepackage{amsmath,amssymb,mathtools,booktabs,xcolor}
\usepackage[protrusion=true,expansion=false]{microtype}
\usepackage[numbers,sort&compress]{natbib}
\usepackage[colorlinks=true,linkcolor=blue!55!black,citecolor=blue!55!black,urlcolor=blue!55!black]{hyperref}
\hypersetup{hypertexnames=false,pdftitle={},pdfauthor={},pdfsubject={},pdfkeywords={},pdfcreator={},pdfproducer={}}
\newtheorem{theorem}{Theorem}[section]
\newtheorem{lemma}[theorem]{Lemma}
\newtheorem{corollary}[theorem]{Corollary}
\theoremstyle{remark}
\newtheorem{remark}[theorem]{Remark}
\newcommand{\calM}{\mathcal M}
\newcommand{\im}{\operatorname{im}}
\title[Least-monomer matching dynamics]{A Least-Monomer Map on Odd-Cycle Matchings\newline and Its Interval Fibres}
\author{Anonymous}
\date{}
\keywords{finite dynamical system, matching, augmenting path, monomer, transient, inverse fibre}
\begin{document}
\begin{abstract}
We study a deterministic map on all matchings of a labelled odd cycle.
When at least three vertices are unmatched, the map flips the alternating
arc from the least unmatched vertex to the next unmatched vertex clockwise.
With one unmatched vertex, it slides the next dimer into that vertex.
Every state enters a single recurrent cycle, and its exact entry time is
its deficiency from maximum matching size. The inverse problem depends on
the labels: if a target has least unmatched vertex $u$, its fibre has size
$\binom{\lfloor u/2\rfloor+1}{2}$, with one additional predecessor when the
target is maximum. We prove this formula by a bijection with intervals of a
forced dimer prefix. It includes all empty fibres, gives the first-image
census, and identifies a unique largest fibre of size
$1+\binom{m+1}{2}$ on a cycle of length $2m+1$. Exhaustive checks through
length twenty-one support reproducibility, not the all-parameter proof.
Classical augmentation and matching enumeration receive no contribution
credit; external owner status remains \texttt{HOLD\_EXTERNAL}.
\end{abstract}
\maketitle

\section{The question and its classical boundary}
An augmenting path explains how to increase the size of a matching.
It does not determine the predecessor sets of a particular deterministic
rule for selecting that path. This note fixes such a rule on an odd cycle
and solves both its recurrent dynamics and its labelled one-step inverse.
The distinction is between finding a maximum matching and describing the
entire finite functional graph of one specified update.

Augmenting-path methods are classical \cite{Berge1957}. The local fact
needed below is proved directly: an alternating path with unmatched
endpoints gains one matching edge when flipped. Matching counts of paths
and cycles, their Fibonacci identities, and the motion of a single defect
in an alternating configuration are also background. None is presented
as a new matching principle. Perfect-matching reconfiguration via
alternating cycles is a related but different problem
\cite{ItoEtAl2022}: there the rank is preserved, whereas the transient
branch here flips a path and strictly raises rank.

The result retained here is the joint object consisting of the
least-monomer scheduler, its odd-cycle recurrent splice, and an exact
target-prefix interval bijection. Internally, hard-core traffic in P90,
the killed path-retile scout GCM, and the stochastic augmentation scout
AP1 remain close comparison surfaces. Their generic alternating-core,
augmentation, and finite-map formulas receive zero credit. Our rule uses
one globally selected arc, and the inverse is a finite interval choice
before the target's least unmatched label. These distinctions do not
constitute a novelty or priority certificate. The bounded owner audit and
its incomplete historical coverage remain visible in the companion record.

\section{The labelled matching map}
Fix $n=2m+1\ge3$. Label the vertices of the simple cycle $C_n$ by
$0,1,\ldots,n-1$ clockwise, and put $e_i=\{i,i+1\pmod n\}$.
Let $\calM_n$ be the set of all matchings, including the empty matching.
A \emph{monomer} is an unmatched vertex. There are $n-2|M|$ monomers in
$M$, an odd positive number. Least labels refer to the ordinary linear
order, not to an unspecified cyclic starting point.

Define $F_n:\calM_n\to\calM_n$ by the following rule. If there are at
least three monomers, let $a$ be the least one and $b$ the next monomer
clockwise from $a$. Flip membership of every edge on the clockwise
$a$--$b$ arc. If $a$ is the only monomer, flip $e_a$ and $e_{a+1}$.
These cases exhaust the carrier and are recomputed from the current state.

\begin{lemma}[closure and rank increase]\label{lem:closure}
The rule is well defined. If $|M|<m$, then $|F_n(M)|=|M|+1$.
If $|M|=m$ and its monomer is $a$, the output has monomer $a+2\pmod n$.
\end{lemma}
\begin{proof}
Between two consecutive monomers every internal vertex must be matched
inside the arc: the only edges leaving its interior meet the unmatched
endpoints. Thus the number of internal vertices is even, the arc length
is odd, and its edge-membership pattern is
$0,1,0,1,\ldots,0$. Flipping gives $1,0,1,0,\ldots,1$, which matches the
endpoints without changing any other matched status. It gains one edge.
With one monomer $a$, deleting $a$ leaves an even path with a unique
perfect matching. In particular $e_a$ is absent and $e_{a+1}$ present.
Their flip replaces the dimer on $a+1,a+2$ by the dimer on $a,a+1$.
The new monomer is $a+2$, including when either edge crosses the label cut.
\end{proof}

\section{The exact entry time and recurrent cycle}
The tail $\tau(M)$ is the least nonnegative iterate at which the orbit of
$M$ reaches a recurrent state.
\begin{theorem}[complete temporal classification]\label{thm:time}
For every $M\in\calM_n$,
\begin{equation}\label{eq:tail}
 \tau(M)=m-|M|.
\end{equation}
The $n$ maximum matchings form one directed cycle of length $n$. There
are no other recurrent states and no fixed points. The largest tail is
$m$, attained only by the empty matching.
\end{theorem}
\begin{proof}
Lemma~\ref{lem:closure} strictly raises rank until rank $m$, so a state of
smaller rank cannot be recurrent and reaches that rank after exactly
$m-|M|$ steps. For each prescribed monomer there is precisely one maximum
matching, by the unique tiling of the even path obtained on deleting that
vertex. On these $n$ states the monomer moves by $+2$. Because
$\gcd(2,n)=1$, this is a single $n$-cycle. It proves both the exact tail
and the recurrent classification. Only rank zero has tail $m$.
\end{proof}

\begin{corollary}[depth census]\label{cor:depth}
The depth generating polynomial is
\begin{equation}\label{eq:depth}
 \sum_{M\in\calM_n}z^{\tau(M)}
 =\sum_{r=0}^{m}\frac{n}{n-r}\binom{n-r}{r}z^{m-r}.
\end{equation}
\end{corollary}
\begin{proof}
For completeness, the classical matching count has a short direct proof.
A path on $v$ vertices has $\binom{v-r}{r}$ matchings of size $r$:
if its selected edge indices are $i_1<\cdots<i_r$ with successive gaps
at least two, subtract $0,1,\ldots,r-1$ to obtain an ordinary $r$-subset
of $v-r$ indices. A cycle matching without the wrap edge therefore has
$\binom{n-r}{r}$ choices. With that edge present, remove its two endpoints
and count $\binom{n-r-1}{r-1}$ choices on the remaining path. At $r=0$
the latter count is zero. Adding gives
$\frac{n}{n-r}\binom{n-r}{r}$, and \eqref{eq:tail} assigns the exponent.
\end{proof}

\section{Every predecessor is a prefix interval}
Fix $Y\in\calM_n$, and let $u$ be its least monomer. There is always a
monomer because $n$ is odd. Put $r=\lfloor u/2\rfloor$ and
$T_r=r(r+1)/2$.

\begin{theorem}[target-resolved inverse]\label{thm:fibre}
Every one-step fibre is given by
\begin{equation}\label{eq:fibre}
 |F_n^{-1}(Y)|=T_{\lfloor u/2\rfloor}+\mathbf1_{\{|Y|=m\}}.
\end{equation}
More precisely, its transient predecessors correspond bijectively to
nonempty consecutive intervals in the $r$ internal dimers before $u$.
When $Y$ is maximum there is one additional, disjoint rotor predecessor.
\end{theorem}
\begin{proof}
All vertices below $u$ are matched. If $u$ is even, the prefix
$0,\ldots,u-1$ is tiled by $e_0,e_2,\ldots,e_{u-2}$. To justify this,
start at $u-1$, which cannot be matched to the monomer $u$, and work
backwards. If $u$ is odd, the same argument forces the internal dimers
$e_1,e_3,\ldots,e_{u-2}$ and then forces vertex zero to use the wrap
edge $e_{n-1}$. In either case there are exactly $r$ consecutive
\emph{internal} prefix dimers; the wrap dimer in the odd case is excluded.

Suppose a nonmaximum source $M$ maps to $Y$, and its chosen endpoint
monomers are $a,b$. The source has at least three monomers. Since $a$ is
the least and $b$ the next, the selected arc does not cross the label cut
and $a<b$. The update removes just these two monomers. Every remaining
monomer is greater than $b$, so $a<b<u$. On the target arc, the flipped
pattern is $1,0,\ldots,1$. Thus that arc consists of a nonempty interval
of the forced prefix dimers just described.

Conversely, choose any such interval. Flip its target edges backwards.
The resulting matching has two new monomers, the endpoints $a<b$ of the
interval, and all former monomers of $Y$. The endpoints satisfy $b<u$.
They are consequently the two least source monomers and are consecutive
clockwise, so the prescribed scheduler selects this very interval and
returns $Y$. Distinct intervals have distinct endpoint pairs and hence
distinct sources. This proves the bijection. There are
$\sum_{j=1}^r j=T_r$ nonempty intervals.

A maximum source stays maximum and has exactly one successor and
predecessor within the rotor. It can therefore contribute to $Y$ exactly
when $Y$ is maximum, in which case its monomer is $u-2\pmod n$.
The transient predecessors have one fewer edge than $Y$, so none is this
rotor predecessor. There are no other branches of the update.
\end{proof}

\begin{corollary}[support and unique extremizer]\label{cor:max}
A nonmaximum target lies in $\im(F_n)$ exactly when $u\ge2$; every
maximum target lies in the image. The unique largest fibre is at the
maximum matching with monomer $n-1$, and its size is
\begin{equation}\label{eq:max}
 1+\frac{m(m+1)}2.
\end{equation}
\end{corollary}
\begin{proof}
The support statement is \eqref{eq:fibre}, since $T_r>0$ exactly when
$r\ge1$. The inequality $u\le2m$ gives $r\le m$. Equality $r=m$
forces $u=2m$, so every other vertex is matched and $Y$ is uniquely the
maximum matching with that monomer. It has the extra rotor predecessor.
All other targets have $r\le m-1$ and a no larger indicator, proving
the strict maximum.
\end{proof}

For example, on $C_7$ the target with edges $e_0,e_2,e_4$ has monomer
six. Its six transient predecessors are the reversals of the six nonempty
intervals of those three dimers. Its seventh predecessor is the maximum
matching with monomer four. Targets with $u=0$ or $1$ have no transient
predecessor at all. These cases show why the rotor indicator cannot be
absorbed into the interval count.

\section{The first image and bounded exact controls}
Write $F_0=0,F_1=1,F_{j+1}=F_j+F_{j-1}$ for Fibonacci numbers and
$L_j=F_{j-1}+F_{j+1}$ for $j\ge1$.
\begin{corollary}[image census]\label{cor:image}
For every odd $n\ge3$,
\begin{equation}\label{eq:image}
 |\im(F_n)|=F_{n-1}+F_{n-3}+2=L_{n-2}+2.
\end{equation}
\end{corollary}
\begin{proof}
First count all targets with $u\ge2$, equivalently those matching both
vertices zero and one. If $e_0$ is present, the remaining path on $n-2$
vertices has $F_{n-1}$ matchings. If $e_0$ is absent, both $e_{n-1}$ and
$e_1$ must be present. For $n\ge5$ their removal leaves a path on $n-4$
vertices, counted by $F_{n-3}$. When $n=3$ this pair is incompatible, so
its count is zero, agreeing with $F_0$. The path count follows from the
recurrence obtained by leaving its first vertex unmatched or matching it
to the next: a path on $v\ge0$ vertices has $F_{v+1}$ matchings.
The two disjoint cases already include all maximum targets with monomer
at least two. Corollary~\ref{cor:max} adds exactly the two maximum
targets with monomer zero or one. The Lucas identity is its definition.
\end{proof}

\begin{table}[ht]
\centering
\caption{Exact complete-carrier checks. These values are consequences of
the proofs; exhaustion is bounded counterexample pressure.}
\begin{tabular}{rrrrrr}
\toprule
$n$ & States & Image & Recurrent & Max tail & Max fibre\\
\midrule
3 & 4 & 3 & 3 & 1 & 2\\
7 & 29 & 13 & 7 & 3 & 7\\
11 & 199 & 78 & 11 & 5 & 16\\
15 & 1364 & 523 & 15 & 7 & 29\\
21 & 24476 & 9351 & 21 & 10 & 56\\
\bottomrule
\end{tabular}
\end{table}

The paper-local standard-library verifier enumerates every matching for
odd $3\le n\le21$, constructs the full transition graph, and computes
tails and cycles by indegree peeling. It compares every target's actual
indegree with \eqref{eq:fibre}, including empty fibres, checks every depth
coefficient, and tests the unique fibre extremizer. A separate endpoint-
set reconstruction checks the complete inverse bijection. Two fresh
processes are compared byte for byte with the retained transcript.
The proofs above, rather than these finite boxes, establish all-size claims.

The scope is restricted to odd simple cycles with the stated numeric
scheduler. Nothing here gives an arbitrary-graph matching algorithm, a
pointwise all-time inverse formula, or an external novelty clearance.
The image census and depth coefficients are classical enumeration applied
to this literal map, not extra independent mechanisms. The companion
source record preserves the P90/GCM/AP1 subtraction and source-access
limitations. The manuscript remains \texttt{OWNER\_AMBER / HOLD\_EXTERNAL}.

\subsection*{Reproducibility and disclosure}
The source, paper-local verifier, frozen stdout, and deterministic build
instructions accompany this internal anonymous note. No empirical or
human-subject data are used. AI-assisted drafting, proof exploration and
exact verification were used; bounded checks and author checks are not
represented as human expert review. Authorship, funding and conflict
declarations are not inferred for this internal draft.
\bibliographystyle{plainnat}
\bibliography{references}
\end{document}
